%%%%     UsersPackagesCommands.tex              %%%%
%%%%     version 2.1.1 du 18/10/2025            %%%%
%%%%     Author: Romain NOËL                    %%%%
%%%%     email: romain.noel@univ-eiffel.fr      %%%%
%%%%--------------------------------------------%%%%

{% raw %}
%%%%---- OS & FONT     -------------------------%%%%
	\usepackage{iftex}
	\usepackage{etoolbox}
	% \usepackage{lmodern} % Loads Latin Modern fonts (which is derived from computer modern (the default font))
		\ifboolexpr{bool {xetex} or bool {luatex}}{
			\DeclareFontShape{T1}{lmr}{bx}{sc}{<-> ssub * cmr/bx/sc}{}
		}{% PDFlatex doesn't require this substitution
		}
	% \usepackage{newcomputermodern} % Loads unicode-math (seems better than modernlatin and computermodern: supports russian, greek, hebrew, old latin AND double square bracket, etc)


%%%%---- TEXT FORMATTING -----------------------%%%%
	\usepackage[shorthands=off,\lang]{babel}
	\usepackage[autostyle]{csquotes} % control quoting style according to babel language
	\usepackage[inline]{enumitem} % improve itemize en enumerate controls
	\usepackage{soul} % wavy underline etc.


	%%% SMALL CAPS (requires graphics for scalebox)
		\ExplSyntaxOn
			% Create a new command to fake Small Caps by scaling uppercase with one mandatory argument.
			\DeclareDocumentCommand{\fakesc}{ m }{
				\textapperance_smallcaps:n { #1 }
			}

			%–– user-adjustable scaling parameters
			\tl_new:N \l__fauxsc_hscale_dim  \tl_set:Nn \l__fauxsc_hscale_dim {.83}
			\tl_new:N \l__fauxsc_vscale_dim  \tl_set:Nn \l__fauxsc_vscale_dim {.72}%.72

			\tl_new:N \l__textapperance_textsmallcaps_input_tl% Initialize string variable
			\cs_new_protected:Npn \textapperance_smallcaps:n #1 {
				% Fully expand contents and store them in the variable
				\tl_set:Nx \l__textapperance_textsmallcaps_input_tl { #1 }

				% Replace everything thats NOT a capital letter NEITHER punctuation by \textsmaller[2]{\uppercase{<match>}}
				\regex_replace_all:nnN { ([^A-Z\ \,\.\;\:\!\?\'\"\(\)\[\]\{\}0-9]+) }{
					% { \c{relscale}{0.75}\cB\{\c{uppercase}\cB\{\0\cE\}\cE\}} }
					\c{scalebox}{
						\c{l__fauxsc_hscale_dim}}[\c{l__fauxsc_vscale_dim}]{
						\cB\{\c{uppercase}\cB\{\0\cE\}\cE\}
					}
				} \l__textapperance_textsmallcaps_input_tl
				\tl_use:N \l__textapperance_textsmallcaps_input_tl % print the result
			}
		\ExplSyntaxOff

	%%%% Personal %%%%
		\NewDocumentCommand{\foreignText}{m}{\emph{#1}} % conventional police for family names
		\makeatletter
			\NewDocumentCommand{\ie}{s}{\IfBooleanTF#1{\foreignText{i.e.\@}}{\foreignText{id est}\ }}
			\NewDocumentCommand{\eg}{s}{\IfBooleanTF#1{\foreignText{e.g.\@}}{\foreignText{exempli gratia}\ }}
			\NewDocumentCommand{\cf}{s}{\IfBooleanTF#1{\foreignText{cf.\@}}{\foreignText{confer}\ }}
			\NewDocumentCommand{\etc}{s}{\IfBooleanTF#1{\foreignText{etc.\@}}{\foreignText{et cetera}\ }}
			\NewDocumentCommand{\etali}{s}{\IfBooleanTF#1{\foreignText{et al.\@}}{\foreignText{et alii}\ }}
			\NewDocumentCommand{\cad}{s}{\IfBooleanTF#1{c.\texthyp{}à\texthyp{}d.\@}{c'est\texthyp{}à\texthyp{}dire\ }}
		\makeatother
		\NewDocumentCommand{\famName}{m}{\fakesc{#1}} %\fakesc conventional police for family names
		\newcommand{\famDash}{\textendash} % conventional dash between two family names
		\newcommand{\famNames}[2]{\famName{#1}\famDash\famName{#2}} \newcommand{\famNamess}[3]{\famName{#1}\famDash\famName{#2}\famDash\famName{#3}}
		\newcommand{\parDash}{\textemdash} % conventional "parenthesis" dash
		\newcommand{\texthyp}{-} % text hyphenation (between two words)
		\newcommand{\md}{\,.} % math period (dot)
		\newcommand{\mc}{\,,} % math comma
		\newcommand{\fd}{.} % float period (dot)


%%%%---- REVIEW --------------------------------%%%%
	% \usepackage{lineno} % to print the number of lines in the pdf
	% \usepackage{comment} % to comment(hide) part of your latex code (NOT needed anymore with the line below) !
		\NewDocumentEnvironment{comment}{+b}{}{}% {comment}{+b}{#1}{} to make it appear
	\usepackage[authormarkuptext=name, commandnameprefix=always]{changes} % to print comments/annotations in your pdf
		%%%[authormarkuptext=id, authormarkup=superscript,highlightmarkup=background]
		\setauthormarkupposition{left}
		%\DeclareRobustCommand{\hlcolor}[2]{{\sethlcolor{#1}\hl{#2}}} % require \usepackage{soul}
		%\sethighlightmarkup{\hlcolor{authorcolor}{#1}}
		\setauthormarkup{\colorbox{authorcolor}{\sffamily\scriptsize\bfseries\textcolor{white}{#1}}\,}
		\sethighlightmarkup{{\IfIsColored{\color{authorcolor}}{}#1}}

		\definecolor{rncolor}{RGB}{31, 140, 65}
		\definechangesauthor[name=Romain, color=rncolor]{RN}


%%%%---- SPACING & JUSTIFICATION ---------------%%%%
%%%%---- LAYOUT & MARGINS ----------------------%%%%
	\usepackage{geometry} % control margins


%%%%---- FOOT PAGE & HEAD ----------------------%%%%
%%%%---- TABLE of CONTENT & FIGURES ------------%%%%
%%%%---- PARTS & CHAPTERS ----------------------%%%%
%%%%---- WATERMARKS ----------------------------%%%%
%%%%---- COVER PAGE ----------------------------%%%%
%%%%---- IMAGES --------------------------------%%%%
	% \usepackage{memoize} %
	% 	\mmzset{
	% 		auto={algorithm}{nomemoize},
	% 		auto={tcolorbox}{memoize},
	% 	}
	\usepackage{graphicx} % standard float figures
	\usepackage{xcolor} % control colors
	\usepackage{placeins} % provide FloatBarrier
	\usepackage{float} % finer control of float and new float env
	\usepackage{subcaption} % Figures multiples
	% \usepackage{epstopdf} % for eps files
	% \usepackage{svg} % for svg files


%%%%---- DRAWING & GRAPHS ----------------------%%%%
	\usepackage{import} % Establish input relative to a directory
	\usepackage{tikz} % LateX framework to "code" drawings
	\usepackage{pgfplots} % Adapt TikZ (& PGF) to plot data/functions
		\pgfplotsset{compat=newest} % set compatibility
		\ifboolexpr{bool {xetex} or bool {luatex}}{
			%\pdfsuppresswarningpagegroup=1 %. no need with lua
		}{
			\pdfsuppresswarningpagegroup=1 % remove some unnecessary warnings
		}


%%%%---- TABLES & BOXES ------------------------%%%%
	% \usepackage{tcolorbox} % colored box with a lot of flexibility
	% 	\tcbuselibrary{skins,theorems} % extra for pseudo(code)
	\usepackage{tabularray} % Typeset tabular and arrays (contains equivalent of longtable, booktabs and dcolumn at least)
		% \UseTblrLibrary{booktabs} % to load extra commands from booktabs
		% \UseTblrLibrary{diagbox} % for cells with a diagbox
		% \UseTblrLibrary{varwidth} % for cell with variable width
		% \UseTblrLibrary{amsmath} % to improve with +array, +matrix, +bmatrix, +Bmatrix, +pmatrix, +vmatrix, +Vmatrix and +cases
		\UseTblrLibrary{siunitx} % for cell using SI formatting
	\usepackage[l3]{csvsimple} % to import CSV file as tables


%%---- SYMBOLS -------------------------------%%%%
	%%%% Check list %%%%%
		\newlist{todolist}{itemize}{2}
		\setlist[todolist]{label=$\circ$}
			\newcommand{\wontdo}{%
				\rlap{$\square$}{%
					\raisebox{-0.25pt}{\textcolor{red}{\LARGE\hspace{0.25pt}\sffamily x}%
					}%
				}%
				\hspace{-0.0pt}%
			}
			\newcommand{\done}{%
				\rlap{$\square$}{\raisebox{1.pt}{\textcolor{green}{\large\hspace{0.25pt}\checkmark}}%
				}%
				\hspace{-2.5pt}%
			}
			\newcommand{\wip}{%
				\rlap{$\square$}{\raisebox{1.65pt}{\scriptsize\hspace{0.75pt}$\sim$}%
				}%
				\hspace{0.9pt}%
			}
			\newcommand{\todos}{$\square$}


%%%%---- MATHS & PHYSICS -----------------------%%%%
	\usepackage{amsmath, amsfonts, amssymb} % American Math Standard
	\usepackage{mathtools} % extension of AMS with tools
	\usepackage{mleftright} % better left and right bracket management
	\usepackage{derivative} % standard and customable derivative operators
	\usepackage{siunitx} % International System units
	% \usepackage{tensorstyles} % display tensors with standards
	% % \usepackage{tensorsop} % work in progress ^^
	% 	\newcommand{\glsarg}{}
	% 	\newcommand{\glsstar}{}
	% 	\NewDocumentCommand{\tsr}{s o m o}{
	% 		\IfBooleanT{#1}{\renewcommand{\glsstar}{*}}
	% 		\IfNoValueTF{#2}{}{\renewcommand{\glsarg}{#2}}
	% 		\IfNoValueTF{#4}{\gls{#3}}{\gls{#3}[#4]}
	% 	}
	\usepackage{keytheorems} % replacement of amsthm, ntheorem, thmtools
		% custom below
	\usepackage{empheq} % emphasize equations
	\usepackage{xsim} % For exercises
		% custom below
	\usepackage[only,llbracket,rrbracket]{stmaryrd} % double square bracket
	% \ifboolexpr{bool {xetex} or bool {luatex}}{
	% 	\usepackage{unicode-math}
	% }{
	% 	% \usepackage[bb=dsserif]{mathalpha}
	% }


%%%%---- CHEMISTRY & BIO -----------------------%%%%
	% \usepackage{chemfig} % draw molecules
	% \usepackage{mhchem} % chemical formulae/equations

%%%%---- CODING --------------------------------%%%%
	\usepackage{listings} % to include piece of codes
		\lstset{
			numbers=left,
			breaklines=true,
			columns=fullflexible,
  			keepspaces=true,
			basicstyle=\ttfamily,
			keywordstyle=\bfseries\color{blue},
			numberstyle=\color{gray}
		}
	\usepackage{pseudo} % for flexible pseudocode
		% \newtcbtheorem{algorithm}{Algorithm}{pseudo/booktabs, float}{alg} % create a floating env for algorithms with tcolorbox
		\floatstyle{ruled} % create a env for algorithms with Float
		\newfloat{algorithm}{htbp}{alg}[section]
		\floatname{algorithm}{Algorithm}


%%%%---- LINKS & REFERENCES ----------------%%%%
	\usepackage{hyperref} % hypertext references (url, mailto, file, jump to links)
		% hyperref: setup below
	%\usepackage[\lang]{varioref} % provide \vref detect and print the page number of float referenced
	\usepackage[sort&compress]{cleveref}% [noabbrev, capitalise] % provide \cref that deduce and print the type of float referenced
		%\makeatletter %
			\crefname{algorithm}{alg.}{algs.}
			% \crefname{tcb@cnt@algorithm}{\textit{Algorithm}}{\textit{Algorithms}}% algo from tcolorbox
			% \crefname{figure}{{Figure}}{{Figures}}
		%\makeatother
	\usepackage{standalone} % to include (and compile) file without their preamble
		\standaloneconfig{
			% group=false,%
			build={
				% build with a special command.
				command={
					\latex\space\latexoptions\space\quote\commandStandaloneBuild\plotsize\string\input{\file}\quote
				}
			}%
		}

		% Define two new length to store the local values of linewidth and baselineskip, because they are computed once at the beginning of the document.
		% Then when called by an external command, they are not updated.
		\newlength{\locallinewidth}
		\newlength{\localbaselineskip}
		% Update the local values.
		\newcommand{\UpdateLineAndBaseLength}{
			\setlength{\localbaselineskip}{\the\baselineskip}
			\setlength{\locallinewidth}{\the\linewidth}
		}
		% Initialize the values.
		\AtBeginDocument{\UpdateLineAndBaseLength{}}

		\makeatletter % new command to pass info to the build standalone.
		\newcommand{\commandStandaloneBuild}{% Warning newline makes differences Here !!!
			\string\newcommand\string{\string\PathS\string}\string{\PathS\string}%
			\string\newcommand\string{\string\PathR\string}\string{\PathR\string}%
			\string\AtBeginDocument\string{%
				\string\setlength\string{\string\textwidth\string}\string{\the\textwidth\string}%
				\string\setlength\string{\string\textheight\string}\string{\the\textheight\string}%
				\string\setlength\string{\string\linewidth\string}\string{\the\locallinewidth\string}%
				\string\renewcommand\string{\string\familydefault\string}\string{\f@family\string}%
				\string\fontencoding\string{\f@encoding\string}%
				\string\fontfamily\string{\f@family\string}%
				\string\fontseries\string{\f@series\string}%
				\string\fontshape\string{\f@shape\string}%
				\string\fontsize\string{\f@size\string}\string{\the\localbaselineskip\string}%
				\string\selectfont%
				\string\setlength\string{\string\baselineskip\string}\string{\the\localbaselineskip\string}%
				\string\stepcounter\string{levelstanda\string}%\string}%
				\string}%
				%\string\reset@font
		}
		\makeatother

		\newcommand{\plotsize}{} % new command to pass pgfplotset to the build standalone
		% redefine the \pgfplotsset for standalone command.
		\newcommand{\UpdatePlotSize}[2]{%
			\renewcommand{\plotsize}{% redefine the \pgfplotsset to pass to the build standalone command.
				\string\AtBeginDocument\string{%
					\string\pgfplotsset\string{#1\string}%
					\string\def\string\svgwidth\string{#2\string}%
					\string}%
			}%
		}

	%\usepackage{hyperref} % CAREFUL when place this usepackage might need to be invoke earlier then other packages
		\hypersetup{
			colorlinks = true,
			allcolors = blue,
			%hidelinks,
		}


%%%%---- INDEX & GLOSSARIES --------------------%%%%
	\usepackage[record, automake, section=section]{glossaries-extra}
		% New glossaries
		\newglossary*{acronym}{Abbreviations}
		\newglossary*{notation}{Notation}
		\newglossary*{symbol}{Nomenclature}

		% new keys must be defined before \GlsXtrLoadResources
		\glsaddstoragekey{unit}{}{\glsentryunit}

		% define style of the abbreviation in the text
		\setabbreviationstyle[acronym]{long-short}
		\setabbreviationstyle[foreignabbreviation]{long-short-user}
		\setabbreviationstyle[inline]{long-short}% apply normal style to the inline class of acronym
		\setabbreviationstyle[footer]{footnote}% apply footnote style to the footer class of acronym

		\GlsXtrLoadResources[
			src={\PathSymbolsEntries}, % entries in symbols.bib
			%selection=all,
			%save-locations=false,% location list not needed
			entry-type-aliases={main=entry, notation=entry, foreignabbreviation=abbreviation},
			field-aliases={
				foreignlong=user1,
				nativelong=long
			},
			type=acronym,
			match={entrytype=(acronym|abbreviation)},
			category={same as original entry}% requires bib2gls v1.4+
			%sort
		]

		\GlsXtrLoadResources[
			src={\PathSymbolsEntries}, % entries in symbols.bib
			%selection=all,
			%save-locations=false,% location list not needed
			entry-type-aliases={main=entry, notation=entry, foreignabbreviation=abbreviation},
			type=symbol,
			match={entrytype=symbol},
			sort-field={category},% sort @symbol entries by category field
			symbol-sort-fallback=description,% sort @symbol entries by name field
			identical-sort-action=name,% sort by description if sort value identical
			sort={letter-nocase},
			% symbol-sort-fallback=name,% sort @symbol entries by name field
			% identical-sort-action=description,% sort by description if sort value identical
		]
		%\glsdisablehyper


%%%%---- BIBLIOGRAPHY --------------------------%%%%
	\usepackage[sort&compress]{natbib}
		\bibliographystyle{plain} %

	% \usepackage[style=alphabetic, % alphabetic, numeric, authoryear
	% 	backend=biber, % bibtex, biber
	% ]{biblatex}
	% 	% Use famName style for the family names.
	% 	\DefineBibliographyExtras{french}{\restorecommand\mkbibnamefamily}
	% 	\renewcommand*{\mkbibnamefamily}[1]{\famName{#1}}

	% 	%% Let us now load the bib files with a custom command similar to natbib.
	% 	%
	% 	% requires etoolbox to loop from a CSV list. (etoolbox is loaded implicitly by ...)
	% 	%\addbibresource[datatype=bibtex]{\Pathbasebib}
	% 	\newcommand{\myListAddBib}[1]{ \addbibresource[datatype=bibtex]{#1} }
	% 	%\forcsvlist{\mylistDo}{\expandafter\Pathbasebib}
	% 	\expandafter\forcsvlist\expandafter\myListAddBib\expandafter{\Pathbasebib}


%%%%---- THEOREM  (settings) ---------------------------%%%%
	% \usepackage{keytheorems} % SETUP

		% Plain style: corollary, lemma, conjecture
		\newkeytheorem{corollary}[name=Corollary, style=plain]
		\newkeytheorem{lemma}[name=Lemma, style=plain]
		\newkeytheorem{conjecture}[name=Conjecture, style=plain]

		% Break style: theorem
		\newkeytheorem{theorem}[name=Theorem, style=plain]
		\newkeytheorem{paradox}[name=Paradox, style=plain]
		%\newkeytheorem{proof}[name=Proof, style=break]

		% Definition style: property, proposition, contraposition, consequence, hypothesis
		\newkeytheorem{property}[name=Property, style=definition]
		\newkeytheorem{proposition}[name=Proposition, style=definition]
		\newkeytheorem{contraposition}[name=Contraposition, style=definition]
		\newkeytheorem{consequence}[name=Consequence, style=definition]
		\newkeytheorem{hypothesis}[name=Hypothesis, style=definition]

		% Defbreak style: Definition, reminder
		\newkeytheorem{definition}[name=Definition, style=definition]
		\newkeytheorem{reminder}[name=Reminder, style=definition]

		% Remark style: remark, example, counterexample, exercise
		\newkeytheorem{remark}[name=Remark, style=remark]
		\newkeytheorem{example}[name=Example, style=remark]
		\newkeytheorem{counterexample}[name=CounterExample, style=remark]
		%\newkeytheorem{exercise}[name=Exercise, style=remark]  % Giving to XSIM


		\definecolor{profOnlyColor}{RGB}{110, 110, 110}
		\newbool{printProfOnly}
		\booltrue{printProfOnly}

	\ExplSyntaxOn
		\NewDocumentEnvironment{myProfOrNotPrintEnv}{+b}{%
			% before myProfOrNotPrintEnv
			\bool_if:NTF \l__mythm_inProfOnly_bool
			{% true
				\ifbool{printProfOnly}{
					\color{profOnlyColor}
					#1
				}{% then to nothing so it will not be printed
					%\begin{comment}
					\begin{keythms_orig_theorem}
				}%
			}{% FALSE
				#1
			}
		}{% after myProfOrNotPrintEnv
			\bool_if:NTF \l__mythm_inProfOnly_bool
			{% true
				\ifbool{printProfOnly}{% nothing special
				}{% eventually close some environments
					%\end{comment}
				}%
			}{% FALSE: nothing to do
			}
		}

		\bool_new:N \l__mythm_inProfOnly_bool
		\bool_set_false:N \l__mythm_inProfOnly_bool

		\keys_define:nn { keytheorems/thmuse }
		{
			profOnly  .choice:,
			profOnly / false .code:n = {
				\bool_set_false:N \l__mythm_inProfOnly_bool
				\bool_set_false:N \l__keythms_thmuse_ignore_bool
			},
			profOnly / true .code:n = {
				\bool_set_true:N \l__mythm_inProfOnly_bool
				\ifbool{printProfOnly}{}{
					\bool_set_true:N \l__keythms_thmuse_ignore_bool
				}
			},
			profOnly  .initial:n = false,
		}
	\ExplSyntaxOff

		\addtotheoremhook{prehead}{\begin{myProfOrNotPrintEnv}}
		\addtotheoremhook{postfoot}{\end{myProfOrNotPrintEnv}}


%%%%---- EXERCISE (XSIM) ---------------------------%%%%
%%%%---- TRANSLATION ---------------------------%%%%
%%%%---- HYPHENATIONS & ACRONYMS --------------%%%%

{% endraw %}
% EoF